\section{Internal Hom}
\label{section-internal-hom}

\noindent
Let $(X, \mathcal{O}_X)$ be a ringed space.
Let $\mathcal{F}$, $\mathcal{G}$ be $\mathcal{O}_X$-modules.
Consider the rule
$$
U \longmapsto \Hom_{\mathcal{O}_X|_U}(\mathcal{F}|_U, \mathcal{G}|_U).
$$
It follows from the discussion in Sheaves, Section
\ref{sheaves-section-glueing-sheaves} that this is a sheaf of
abelian groups. In addition, given an element
$\varphi \in \Hom_{\mathcal{O}_X|_U}(\mathcal{F}|_U, \mathcal{G}|_U)$
and a section $f \in \mathcal{O}_X(U)$ then we can define
$f\varphi \in \Hom_{\mathcal{O}_X|_U}(\mathcal{F}|_U, \mathcal{G}|_U)$
by either precomposing with multiplication by $f$ on $\mathcal{F}|_U$
or postcomposing with multiplication by $f$ on $\mathcal{G}|_U$ (it gives
the same result). Hence we in fact get a sheaf of $\mathcal{O}_X$-modules.
We will denote this sheaf
$\SheafHom_{\mathcal{O}_X}(\mathcal{F}, \mathcal{G})$.
There is a canonical ``evaluation'' morphism
$$
\mathcal{F}
\otimes_{\mathcal{O}_X}
\SheafHom_{\mathcal{O}_X}(\mathcal{F}, \mathcal{G})
\longrightarrow
\mathcal{G}.
$$
For every $x \in X$ there is also a canonical morphism
$$
\SheafHom_{\mathcal{O}_X}(\mathcal{F}, \mathcal{G})_x
\to
\Hom_{\mathcal{O}_{X, x}}(\mathcal{F}_x, \mathcal{G}_x)
$$
which is rarely an isomorphism.

\begin{lemma}
\label{lemma-internal-hom}
Let $(X, \mathcal{O}_X)$ be a ringed space.
Let $\mathcal{F}$, $\mathcal{G}$, $\mathcal{H}$ be $\mathcal{O}_X$-modules.
There is a canonical isomorphism
$$
\SheafHom_{\mathcal{O}_X}
(\mathcal{F} \otimes_{\mathcal{O}_X} \mathcal{G}, \mathcal{H})
\longrightarrow
\SheafHom_{\mathcal{O}_X}
(\mathcal{F}, \SheafHom_{\mathcal{O}_X}(\mathcal{G}, \mathcal{H}))
$$
which is functorial in all three entries (sheaf Hom in
all three spots). In particular, to give a
morphism $\mathcal{F} \otimes_{\mathcal{O}_X} \mathcal{G} \to \mathcal{H}$
is the same as giving a morphism
$\mathcal{F} \to \SheafHom_{\mathcal{O}_X}(\mathcal{G}, \mathcal{H})$.
\end{lemma}

\begin{proof}
This is the analogue of
Algebra, Lemma \ref{algebra-lemma-hom-from-tensor-product}.
The proof is the same, and is omitted.
\end{proof}

\begin{lemma}
\label{lemma-internal-hom-exact}
Let $(X, \mathcal{O}_X)$ be a ringed space.
Let $\mathcal{F}$, $\mathcal{G}$ be $\mathcal{O}_X$-modules.
\begin{enumerate}
\item If $\mathcal{F}_2 \to \mathcal{F}_1 \to \mathcal{F} \to 0$
is an exact sequence of $\mathcal{O}_X$-modules, then
$$
0 \to
\SheafHom_{\mathcal{O}_X}(\mathcal{F}, \mathcal{G}) \to
\SheafHom_{\mathcal{O}_X}(\mathcal{F}_1, \mathcal{G}) \to
\SheafHom_{\mathcal{O}_X}(\mathcal{F}_2, \mathcal{G})
$$
is exact.
\item If $0 \to \mathcal{G} \to \mathcal{G}_1 \to \mathcal{G}_2$
is an exact sequence of $\mathcal{O}_X$-modules, then
$$
0 \to
\SheafHom_{\mathcal{O}_X}(\mathcal{F}, \mathcal{G}) \to
\SheafHom_{\mathcal{O}_X}(\mathcal{F}, \mathcal{G}_1) \to
\SheafHom_{\mathcal{O}_X}(\mathcal{F}, \mathcal{G}_2)
$$
is exact.
\end{enumerate}
\end{lemma}

\begin{proof}
Let $\mathcal{F}_2 \to \mathcal{F}_1 \to \mathcal{F} \to 0$
be as in (1). For every $U \subset X$ open the sequence
$$
0 \to \Hom_{\mathcal{O}_U}(\mathcal{F}|_U, \mathcal{G}|_U) \to
\Hom_{\mathcal{O}_U}(\mathcal{F}_1|_U, \mathcal{G}|_U) \to
\Hom_{\mathcal{O}_U}(\mathcal{F}_2|_U, \mathcal{G}|_U)
$$
is exact by Homology, Lemma \ref{homology-lemma-check-exactness}.
This means that
taking sections over $U$ of the sequence of sheaves in (1) produces
an exact sequence of abelian groups. Hence the sequence in (1) is exact
by definition. The proof of (2) is exactly the same.
\end{proof}

\begin{lemma}
\label{lemma-adjoint-tensor-restrict}
Let $X$ be a topological space. Let $\mathcal{O}_1 \to \mathcal{O}_2$
be a homomorphism of sheaves of rings. Then we have
$$
\Hom_{\mathcal{O}_1}(\mathcal{F}_{\mathcal{O}_1}, \mathcal{G}) =
\Hom_{\mathcal{O}_2}(\mathcal{F},
\SheafHom_{\mathcal{O}_1}(\mathcal{O}_2, \mathcal{G}))
$$
bifunctorially in $\mathcal{F} \in \textit{Mod}(\mathcal{O}_2)$
and $\mathcal{G} \in \textit{Mod}(\mathcal{O}_1)$.
\end{lemma}

\begin{proof}
Omitted. This is the analogue of
Algebra, Lemma \ref{algebra-lemma-adjoint-hom-restrict}
and is proved in exactly the same way.
\end{proof}

\begin{lemma}
\label{lemma-stalk-internal-hom}
Let $(X, \mathcal{O}_X)$ be a ringed space.
Let $\mathcal{F}$, $\mathcal{G}$ be $\mathcal{O}_X$-modules.
If $\mathcal{F}$ is of finite type then the canonical map
$$
\SheafHom_{\mathcal{O}_X}(\mathcal{F}, \mathcal{G})_x
\to
\Hom_{\mathcal{O}_{X, x}}(\mathcal{F}_x, \mathcal{G}_x)
$$
is injective.
If $\mathcal{F}$ is finitely presented,
this canonical morphism is an isomorphism.
\end{lemma}

\begin{proof}
The map sends the equivalence class of
$(U, \varphi)$ in $\SheafHom_{\mathcal{O}_X}(\mathcal{F}, \mathcal{G})_x$,
where $x \in U \subset X$ is open and
$\varphi \in \Hom_{\mathcal{O}_U}(\mathcal{F}|_U, \mathcal{G}|_U)$,
to the induced map on stalks at $x$, namely
$\varphi_x : \mathcal{F}_x \to \mathcal{G}_x$.

\medskip\noindent
Suppose $\mathcal{F}$ is of finite type.
Pick a representative $(U, \varphi)$ of an element $\sigma$
in the kernel of the map, i.e., $\varphi_x = 0$.
Shrinking $U$ if necessary, choose sections
$s^1, \ldots, s^n \in \mathcal{F}(U)$ generating $\mathcal{F}|_U$.
Since $\varphi_x(s^i_x) = 0$
and we are dealing with a finite number of sections,
we can find an open neighborhood 
$V \subset U$ of $x$ such that $\varphi_V(s^i|_V)=0$
for all $i = 1, \ldots, n$. Since $s^i|_V$, $i = 1, \ldots, n$ generate
$\mathcal{F}|_V$ this means that $\varphi|_V = 0$.
Since $(U, \varphi)$ is equivalent to $(V, \varphi|_V)$
we conclude $\sigma = 0$ and injectivity of the map follows.

\medskip\noindent
Next, assume $\mathcal{F}$ is finitely presented.
By localizing on $X$ we may assume that $\mathcal{F}$ has a presentation
$$
\bigoplus\nolimits_{j = 1, \ldots, m}
\mathcal{O}_X
\longrightarrow
\bigoplus\nolimits_{i = 1, \ldots, n}
\mathcal{O}_X
\to
\mathcal{F}
\to
0.
$$
By Lemma \ref{lemma-internal-hom-exact} this gives an exact sequence
$
0 \to
\SheafHom_{\mathcal{O}_X}(\mathcal{F}, \mathcal{G}) \to
\bigoplus\nolimits_{i = 1, \ldots, n} \mathcal{G}
\longrightarrow
\bigoplus\nolimits_{j = 1, \ldots, m} \mathcal{G}.
$
Taking stalks we get an exact sequence
$
0 \to
\SheafHom_{\mathcal{O}_X}(\mathcal{F}, \mathcal{G})_x \to
\bigoplus\nolimits_{i = 1, \ldots, n} \mathcal{G}_x
\longrightarrow
\bigoplus\nolimits_{j = 1, \ldots, m} \mathcal{G}_x
$
and the result follows since $\mathcal{F}_x$ sits in
an exact sequence
$
\bigoplus\nolimits_{j = 1, \ldots, m}
\mathcal{O}_{X, x}
\longrightarrow
\bigoplus\nolimits_{i = 1, \ldots, n}
\mathcal{O}_{X, x}
\to
\mathcal{F}_x
\to
0
$
which induces the exact sequence
$
0 \to
\Hom_{\mathcal{O}_{X, x}}(\mathcal{F}_x, \mathcal{G}_x) \to
\bigoplus\nolimits_{i = 1, \ldots, n} \mathcal{G}_x
\longrightarrow
\bigoplus\nolimits_{j = 1, \ldots, m} \mathcal{G}_x
$
which is the same as the one above.
\end{proof}

\begin{lemma}
\label{lemma-pullback-internal-hom}
Let $f : (X, \mathcal{O}_X) \to (Y, \mathcal{O}_Y)$ be a morphism
of ringed spaces. Let $\mathcal{F}$, $\mathcal{G}$ be $\mathcal{O}_Y$-modules.
If $\mathcal{F}$ is finitely presented and $f$ is flat,
then the canonical map
$$
f^*\SheafHom_{\mathcal{O}_Y}(\mathcal{F}, \mathcal{G})
\longrightarrow
\SheafHom_{\mathcal{O}_X}(f^*\mathcal{F}, f^*\mathcal{G})
$$
is an isomorphism.
\end{lemma}

\begin{proof}
Note that $f^*\mathcal{F}$ is also finitely presented
(Lemma \ref{lemma-pullback-finite-presentation}).
Let $x \in X$ map to $y \in Y$. Looking at the stalks
at $x$ we get an isomorphism by
Lemma \ref{lemma-stalk-internal-hom} and
More on Algebra, Lemma
\ref{more-algebra-lemma-pseudo-coherence-and-base-change-ext}
to see that in this case $\Hom$ commutes with base change by
$\mathcal{O}_{Y, y} \to \mathcal{O}_{X, x}$.
Second proof: use the exact same argument as given
in the proof of Lemma \ref{lemma-stalk-internal-hom}.
\end{proof}

\begin{lemma}
\label{lemma-internal-hom-locally-kernel-direct-sum}
Let $(X, \mathcal{O}_X)$ be a ringed space.
Let $\mathcal{F}$, $\mathcal{G}$ be $\mathcal{O}_X$-modules.
If $\mathcal{F}$ is finitely presented then the sheaf
$\SheafHom_{\mathcal{O}_X}(\mathcal{F}, \mathcal{G})$ is
locally a kernel of a map between finite direct sums
of copies of $\mathcal{G}$.
In particular, if $\mathcal{G}$ is coherent then
$\SheafHom_{\mathcal{O}_X}(\mathcal{F}, \mathcal{G})$
is coherent too.
\end{lemma}

\begin{proof}
The first assertion
we saw in the proof of Lemma \ref{lemma-stalk-internal-hom}.
And the result for coherent sheaves then follows from
Lemma \ref{lemma-coherent-abelian}.
\end{proof}

\begin{lemma}
\label{lemma-hom-finite-presentation-colimit}
Let $X$ be a ringed space. Let $\mathcal{F}$ be an $\mathcal{O}_X$-module
of finite presentation. Let
$\mathcal{G} = \colim_{\lambda \in \Lambda} \mathcal{G}_\lambda$
be a filtered colimit of $\mathcal{O}_X$-modules. Then the canonical map
$$
\colim_\lambda \SheafHom_{\mathcal{O}_X}(\mathcal{F}, \mathcal{G}_\lambda)
\longrightarrow
\SheafHom_{\mathcal{O}_X}(\mathcal{F}, \mathcal{G})
$$
is an isomorphism.
\end{lemma}

\begin{proof}
Taking colimits of sheaves of modules commutes with restriction to opens, see
Sheaves, Section \ref{sheaves-section-limits-sheaves}. Hence we may assume
$\mathcal{F}$ has a global presentation
$$
\bigoplus\nolimits_{j = 1, \ldots, m}
\mathcal{O}_X
\longrightarrow
\bigoplus\nolimits_{i = 1, \ldots, n}
\mathcal{O}_X
\to
\mathcal{F}
\to
0
$$
The functor $\SheafHom_{\mathcal{O}_X}(-, -)$ commutes with finite direct sums
in either variable and $\SheafHom_{\mathcal{O}_X}(\mathcal{O}_X, -)$
is the identity functor. By this and by Lemma \ref{lemma-internal-hom-exact}
we obtain an exact sequence
$$
0 \to
\SheafHom_{\mathcal{O}_X}(\mathcal{F}, \mathcal{G}) \to
\bigoplus\nolimits_{i = 1, \ldots, n} \mathcal{G} \to
\bigoplus\nolimits_{j = 1, \ldots, m} \mathcal{G}
$$
Since filtered colimits are exact in $\textit{Mod}(\mathcal{O}_X)$
also the top row in the following commutative diagram is exact
$$
\xymatrix{
0 \ar[r] &
\colim_\lambda \SheafHom_{\mathcal{O}_X}(\mathcal{F}, \mathcal{G}_\lambda)
\ar[r] \ar[d] &
\colim_\lambda \bigoplus\nolimits_{i = 1, \ldots, n} \mathcal{G}_\lambda
\ar[r] \ar[d] &
\colim_\lambda \bigoplus\nolimits_{j = 1, \ldots, m} \mathcal{G}_\lambda
\ar[d] \\
0 \ar[r] &
\SheafHom_{\mathcal{O}_X}(\mathcal{F}, \mathcal{G}) \ar[r] &
\bigoplus\nolimits_{i = 1, \ldots, n} \mathcal{G} \ar[r] &
\bigoplus\nolimits_{j = 1, \ldots, m} \mathcal{G}
}
$$
Since the right two vertical arrows are isomorphisms we conclude.
\end{proof}

\begin{lemma}
\label{lemma-finite-presentation-quasi-compact-colimit}
Let $X$ be a ringed space.
Let $I$ be a preordered set and
let $(\mathcal{F}_i, \varphi_{ii'})$ be a system over $I$
consisting of sheaves of $\mathcal{O}_X$-modules
(see Categories, Section \ref{categories-section-posets-limits}).
Assume
\begin{enumerate}
\item $I$ is directed,
\item $\mathcal{G}$ is an $\mathcal{O}_X$-module of finite presentation, and
\item $X$ has a cofinal system of open coverings
$\mathcal{U} : X = \bigcup_{j\in J} U_j$ with
$J$ finite and $U_j \cap U_{j'}$ quasi-compact
for all $j, j' \in J$.
\end{enumerate}
Then we have
$$
\colim_i \Hom_X(\mathcal{G}, \mathcal{F}_i)
=
\Hom_X(\mathcal{G}, \colim_i \mathcal{F}_i).
$$
\end{lemma}

\begin{proof}
Set
$\mathcal{H} = \SheafHom_{\mathcal{O}_X}(\mathcal{G}, \colim \mathcal{F}_i)$
and $\mathcal{H}_i = \SheafHom_{\mathcal{O}_X}(\mathcal{G}, \mathcal{F}_i)$.
Recall that
$$
\Hom_X(\mathcal{G}, \mathcal{F}) = \Gamma(X, \mathcal{H})
\quad\text{and}\quad
\Hom_X(\mathcal{G}, \mathcal{F}_i) = \Gamma(X, \mathcal{H}_i)
$$
by construction. By Lemma \ref{lemma-hom-finite-presentation-colimit} we have
$\mathcal{H} = \colim \mathcal{H}_i$. Thus the lemma follows from
Sheaves, Lemma \ref{sheaves-lemma-directed-colimits-sections}.
\end{proof}

\begin{remark}
\label{remark-condition-necessary}
In the lemma above some condition beyond the condition that $X$
is quasi-compact is necessary. See
Sheaves, Example \ref{sheaves-example-conditions-needed-colimit}.
\end{remark}




